% =============================================================
%  Module Description — Introduction to Digital Design and Computer Architecture (CS102-F26)
%  Shown on https://wosm.academy as the course page for CS102-F26.
%  Compile with: pdflatex cs102-module-description.tex
%  Write "TODO(owner): ..." where information is still missing;
%  the website highlights it.
% =============================================================
\documentclass[11pt,a4paper]{article}

\usepackage[utf8]{inputenc}
\usepackage[T1]{fontenc}
\usepackage{lmodern}
\usepackage[margin=2.2cm]{geometry}
\usepackage{booktabs}
\usepackage{tabularx}
\usepackage{array}
\usepackage{enumitem}
\usepackage{xcolor}
\usepackage{titlesec}
\usepackage{fancyhdr}
\usepackage[hidelinks]{hyperref}

% ---------- Module details (edit these) ----------------------
% \ModuleCode must equal the course's `code` in src/content/courses/.
\newcommand{\ModuleTitle}{Introduction to Digital Design and Computer Architecture}
\newcommand{\ModuleCode}{CS102-F26}
\newcommand{\Programme}{Foundation programme}
\newcommand{\Level}{Semester 1}
\newcommand{\Semester}{Fall 2026}
\newcommand{\Credits}{TODO(owner)}
\newcommand{\ModuleLeader}{Dr. Ali Khudiyev}
\newcommand{\Email}{TODO(owner)}
\newcommand{\OfficeHours}{TODO(owner)}
\newcommand{\Prerequisites}{None}

% ---------- Styling ------------------------------------------
\definecolor{accent}{RGB}{0,70,130}
\titleformat{\section}{\large\bfseries\color{accent}}{\thesection.}{0.5em}{}
\titlespacing*{\section}{0pt}{1.2em}{0.5em}
\setlist{itemsep=2pt, topsep=4pt}
\setlength{\parindent}{0pt}
\setlength{\parskip}{4pt}

\pagestyle{fancy}
\fancyhf{}
\lhead{\small \ModuleCode\ -- \ModuleTitle}
\rhead{\small Module Description}
\cfoot{\small \thepage}
\renewcommand{\headrulewidth}{0.4pt}

% =============================================================
\begin{document}

\begin{center}
  {\LARGE\bfseries\color{accent} \ModuleTitle}\\[4pt]
  {\large Module Description}
\end{center}
\vspace{6pt}

% ---------- Summary table (not shown on the website: the course page
% shows these details from the macros above) --------------------------
\begin{tabularx}{\textwidth}{>{\bfseries}l X >{\bfseries}l X}
\toprule
Module code   & \ModuleCode    & Credits      & \Credits \\
Programme     & \Programme     & Level        & \Level \\
Semester      & \Semester      & Prerequisites & \Prerequisites \\
Module leader & \ModuleLeader  & Email        & \href{mailto:\Email}{\Email} \\
Office hours  & \multicolumn{3}{l}{\OfficeHours} \\
\bottomrule
\end{tabularx}

% ---------- Sections (shown on the website) ------------------
\section{Module Overview}
This course is for people who are curious about inner working principles of
computers. You will learn about computer systems (e.g.\ basic architecture and
organization, Von-Neumann-Computer, machine-instruction cycle, hardware-software
interface), instruction set architecture (ISA), fundamental methods of circuit
design, structural organization of processor using circuit elements, processor
design variants and their implications, some advanced aspects of description and
implementation and verification of digital circuits, and other things such as
microprocessor architectures and systems, parallel and distributed systems, memory
systems, and input/output (I/O) systems.

\section{Aims}
\begin{itemize}
  \item TODO(owner): aims of the module.
\end{itemize}

\section{Learning Outcomes}
After completing this course, you will be able to:
\begin{enumerate}[label=LO\arabic*.]
  \item Understand computer systems and layered abstract machines.
  \item Gain an introductory overview of computer architecture and fundamental
        knowledge of low-level programming, digital circuits, and processor design.
  \item Apply basic concepts of computer architecture and explain machine
        instruction cycles based on the underlying hardware at the
        register-transfer level.
  \item Classify different types of computer architectures and understand the
        basic principles of modern computer architecture.
\end{enumerate}

\section{Syllabus}
\begin{tabularx}{\textwidth}{c X}
\toprule
\textbf{Unit} & \textbf{Topic} \\
\midrule
1  & Introduction \& Data Representation \\
2  & Assembly Languages \\
3  & Von Neumann Architecture \\
4  & Other ISAs \\
5  & System Architecture \\
6  & Memory Management \& Cache \\
7  & Boolean Algebra \\
8  & Adders \\
9  & Multiplication \\
10 & Memory \\
11 & Processor Architecture \\
12 & Finite Automata \\
13 & Multi-cycle Processor \\
14 & Modern Systems \& Parallelism \\
15 & The I/O System \\
16 & Physical Layer Implementation \\
17 & Pipelining \\
18 & K-Maps \\
19 & Binary Decision Diagrams \\
20 & SAT \\
21 & Physical Design \\
22 & Review \\
\bottomrule
\end{tabularx}

\section{Teaching and Learning Methods}
The module is self-paced. All materials are provided on Moodle.

\begin{tabularx}{0.6\textwidth}{X r}
\toprule
\textbf{Activity} & \textbf{Hours} \\
\midrule
Lectures            & TODO(owner) \\
Labs / Tutorials    & TODO(owner) \\
Independent study   & TODO(owner) \\
\midrule
\textbf{Total}      & \textbf{TODO(owner)} \\
\bottomrule
\end{tabularx}

\section{Assessment}
Your understanding and knowledge of the concepts presented will be assessed
through quizzes, projects, and challenges. Any permitted support materials will be
provided during the examination; no additional aids or resources are allowed.

TODO(owner): weight of each assessment component and the learning outcomes it covers.

\section{Reading List}
\begin{itemize}
  \item David A.\ Patterson, John L.\ Hennessy. \textit{Computer Organization and
        Design: The Hardware/Software Interface}.
  \item Sarah L.\ Harris, David Harris. 2022. \textit{Digital Design and Computer
        Architecture: RISC-V Edition}. Elsevier.
  \item Andrew B.\ Kahng, Jens Lienig, Igor L.\ Markov, Jin Hu. 2022. \textit{VLSI
        Physical Design: From Graph Partitioning to Timing Closure}. Springer
        International Publishing.
  \item \href{https://docs.riscv.org/reference/isa/_attachments/riscv-privileged.pdf}{RISC-V ISA}
\end{itemize}

\section{Materials and Licence}
Course materials are published under the
\href{https://creativecommons.org/licenses/by-sa/4.0/}{Creative Commons Attribution-ShareAlike 4.0}
licence unless stated otherwise. See \href{https://wosm.academy/licence/}{wosm.academy/licence} for details.

\section{Policies}
\begin{itemize}
  \item \textbf{Attendance:} TODO(owner)
  \item \textbf{Late submissions:} TODO(owner)
  \item \textbf{Academic integrity:} All submitted work must be your own.
        See the \href{https://wosm.academy/code-of-conduct/}{Code of Conduct}.
\end{itemize}

\end{document}
